\chapter{Esperimenti e risultati}
\label{ch:risultati}

% BOZZA — per ora contiene gli esperimenti preliminari (appunti, agosto
% 2026). Il setup sperimentale definitivo (benchmark, metriche, baseline)
% va concordato con la relatrice.

\section{Setup sperimentale}
% Definito con la relatrice (4 agosto 2026):
% - dataset: UnLeakedTestBench, ULT_Lite.jsonl (20 campioni, obiettivo 100).
%   Ogni campione: func_name, code (funzione autonoma), prompt (descrizione
%   in linguaggio naturale), task_id, test_list (assert di riferimento umani).
% - modelli: tre Llama di taglia crescente (1B, 3B, 8B) dal catalogo NVIDIA.
% - condizioni: temperatura 0, template di prompt fissato e dichiarato,
%   output salvati in file separati.
% - baseline secondaria: generazione senza LLM (AST + Klara).

\section{Metriche}
% - coverage del codice sotto test;
% - numero di esecuzioni per riga (quante volte la stessa riga viene rieseguita);
% - righe duplicate tra i file di test (duplicate-code-detection-tool);
% - classificazione degli esiti: passato / eseguibile ma fallito / non eseguibile.
% TODO: spiegazione visiva delle metriche + esempio end-to-end di una generazione.

\section{Esperimenti preliminari}

\subsection{Validazione della meccanica del ciclo}
% Esperimento manuale (LLM via chat, copia-incolla umano come
% orchestratore): test per una sola funzione -> coverage 47% ->
% feedback con le righe mancanti -> coverage 100% in una iterazione.
% Nota critica: su funzioni semplici il 100% è quasi garantito; l'esperimento
% valida la meccanica del ciclo, non la sua utilità.

\subsection{Funzioni complesse e codice irraggiungibile}
% Bersaglio con rami annidati, eccezioni e casi limite (49 statement):
% 46 test generati al primo giro, tutti passati, coverage 98%.
% L'unica riga scoperta è dimostrabilmente irraggiungibile (codice
% difensivo morto): nessun test può coprirla. Implicazioni: criterio di
% stop del loop e valore diagnostico delle righe residue.

\section{Risultati}
% TODO: tabelle e grafici (coverage per configurazione e per iterazione)
% quando il sistema automatico sarà operativo.
